\documentclass[a4paper,12pt]{article}
\usepackage[T2A]{fontenc}
\usepackage[utf8]{inputenc}
\usepackage[russian]{babel}
\usepackage{amsmath,amssymb,mathtools}
\usepackage{PTSans}
\renewcommand{\familydefault}{\sfdefault}
\usepackage{titlesec}
\titleformat{\subsection}{\normalfont\bfseries\normalsize\color{ink}}{}{0pt}{}
\titlespacing*{\subsection}{0pt}{1.5ex plus .3ex}{.3ex}
\usepackage{geometry}
\geometry{left=23mm,right=23mm,top=20mm,bottom=22mm}
\usepackage{microtype}
\usepackage{tabularx,booktabs,array}
\usepackage{xcolor}
\usepackage{tikz}
\usepackage{hyperref}
\definecolor{accent}{HTML}{006E70}
\definecolor{ink}{HTML}{252D32}
\definecolor{muted}{HTML}{606B70}
\definecolor{critical}{HTML}{A12C31}
\hypersetup{colorlinks=true,linkcolor=accent,pdfauthor={Лёвин Артём Александрович},pdftitle={Ревью домашней работы — Анна Баннова}}
\setlength{\parindent}{0pt}
\setlength{\parskip}{4pt}
\renewcommand{\arraystretch}{1.3}
\pagestyle{plain}
\begin{document}

\pagenumbering{roman}
\begin{titlepage}
\centering
\vspace*{21mm}
{\large\color{muted} МАТЕМАТИКА\par}
\vspace{20mm}
{\Huge\bfseries\color{ink} Ревью домашней\\[5pt] работы\par}
\vspace{13mm}
{\large Задачи на встречное движение\par}
\vspace{19mm}
{\Large Анна Баннова\par}
\vspace{8mm}
{\large 5 ноября 2026 г.\par}
\vfill
{\color{muted}\normalsize Лёвин Артём Александрович\\[4pt]
эксклюзивно для Анны Банновой\par}
\vspace{13mm}
\begin{tikzpicture}
\draw[accent!75,line width=.9pt] (-4,0) .. controls (-2.1,.32) and (-.7,-.23) .. (1.1,0) .. controls (2.4,.19) and (3.1,.14) .. (4,0);
\end{tikzpicture}
\end{titlepage}
\pagenumbering{arabic}

\section*{Сводная таблица}
\vspace{3mm}
\small
\begin{tabularx}{\linewidth}{@{}>{\raggedright\arraybackslash}p{13mm}>{\raggedright\arraybackslash}p{27mm}X>{\raggedright\arraybackslash}p{21mm}>{\raggedright\arraybackslash}p{24mm}@{}}
\toprule
\textbf{№} & \textbf{Тема} & \textbf{Суть замечания} & \textbf{Ваш ответ} & \textbf{Верный ответ}\\
\midrule
\hyperref[task:10]{10} & Встречное движение с задержкой старта & В последней строке первая цифра промежуточного числа похожа на 9; для равенства требуется 4. Итог найден верно. & $297{,}5$ км & $297{,}5$ км\\
\bottomrule
\end{tabularx}
\normalsize
\vspace{8mm}
\textcolor{muted}{\small В таблицу включён только участок решения, требующий уточнения.}

\newpage
\thispagestyle{empty}
\vspace*{\fill}
\begin{center}
{\Huge\bfseries\color{ink} Разбор ошибок}
\end{center}
\vspace*{\fill}
\newpage

\phantomsection\label{task:10}
\section*{Задание № 10}
\subsection*{Текст задания}
Между городами А и В 655 км. Первый автомобиль выехал из А со скоростью 65 км/ч, а через 2 часа второй выехал навстречу из В со скоростью 85 км/ч. Найдите расстояние от В до места встречи. Ответ запишите в километрах.

\subsection*{Ваше решение}
Вы обозначили через $x$ время движения первого автомобиля до встречи. Второй автомобиль ехал на 2 часа меньше. Получилось правильное уравнение:
\[
65x+85(x-2)=655.
\]
Далее верно выполнены преобразования:
\[
65x+85x-170=655,
\qquad 150x=825,
\qquad x=5{,}5.
\]
Для расстояния от В Вы использовали выражение $85(5{,}5-2)$.

\subsection*{Ваш ответ}
$297{,}5$ км. \textbf{Он верный.}

\subsection*{Ошибка в оформлении вычисления}
\textcolor{critical}{\textbf{Ошибка:}} в заключительной строке промежуточное число на фотографии выглядит как $967{,}5$. При раскрытии скобок должно получиться $467{,}5$. Возможно, это особенность написания цифры 4; однако такая запись допускает неверное прочтение.

\subsection*{Где именно возникает неточность}
Последний полностью правильный шаг:
\[
S_{\text{В}}=85(5{,}5-2).
\]
Запись после раскрытия скобок следует оформить однозначно:
\[
85\cdot5{,}5-85\cdot2
=\underbrace{467{,}5}_{85\cdot5{,}5}-170.
\]
Если первое число прочитать как $967{,}5$, то равенство с $297{,}5$ уже неверно: $967{,}5-170=797{,}5$.

\subsection*{Почему здесь нужно быть внимательной}
При умножении разности на число множитель умножается \emph{на каждое} слагаемое:
\[
a(b-c)=ab-ac.
\]
Поэтому $85(5{,}5-2)=85\cdot5{,}5-85\cdot2$. Необходимо отдельно проверить оба произведения и только затем выполнить вычитание. Поскольку окончательное расстояние у Вас верное, замечание относится к читаемости промежуточной арифметической записи, а не к выбору метода.

\subsection*{\textcolor{accent}{Ключевая идея}}
Если второй участник начал движение позже на $2$ часа, то при общем времени $x$ от первого старта он ехал $x-2$ часа. До встречи сумма пройденных расстояний равна расстоянию между городами.

\subsection*{Исправление}
Сохраняем Ваше правильное значение $x=5{,}5$ и переписываем только последний расчёт:
\[
S_{\text{В}}=85(5{,}5-2)=85\cdot3{,}5=297{,}5\text{ км}.
\]

\subsection*{Полное правильное решение}
Пусть $x$ --- число часов от выезда первого автомобиля до встречи. Второй выехал через 2 часа, поэтому его время в пути $x-2$ часов. Поскольку второй успел выехать до встречи, $x>2$.

Первый автомобиль прошёл $65x$ километров, второй --- $85(x-2)$ километров. Они двигались навстречу, значит сумма пройденных расстояний равна $655$ км:
\[
\begin{aligned}
65x+85(x-2)&=655,\\
65x+85x-170&=655,\\
150x&=825,\\
x&=5{,}5.
\end{aligned}
\]
Второй автомобиль находился в дороге $5{,}5-2=3{,}5$ часа. Значит, от города В до места встречи он проехал
\[
85\cdot3{,}5=297{,}5\text{ км}.
\]

\subsection*{Проверка}
Первый автомобиль прошёл $65\cdot5{,}5=357{,}5$ км. Сумма расстояний:
\[
357{,}5+297{,}5=655\text{ км}.
\]
Она совпадает с расстоянием между городами; условие $x>2$ соблюдено.

\subsection*{\textcolor{accent}{Итог}}
Правильный ответ --- \textbf{$297{,}5$ км}. Уравнение, время встречи и конечный результат найдены верно. При записи промежуточных чисел проследите, чтобы цифры 4 и 9 различались однозначно.

\subsection*{Задача на закрепление}
Между городами А и В 660 км. Первый автомобиль выехал из А со скоростью 75 км/ч. Через 2 часа навстречу ему из В выехал второй автомобиль со скоростью 95 км/ч. На каком расстоянии от В произошла встреча?

\vfill
\begin{center}\small\color{muted}Лёвин Артём Александрович эксклюзивно для Анны Банновой\end{center}
\end{document}
